\section{Exact complementary separation}

We now exhibit the language collision that forces a large inverse window.  Let
the notation of \cref{lem:alternating} stand, and put
\begin{equation}\label{eq:PQ-def}
 Q=(01)^m,\qquad P=(10)^m,\qquad m=2^N-1.
\end{equation}
Then $|P|=|Q|=L-1$ and
\begin{equation}\label{eq:PQ-identities}
 U=Q0=0P,\qquad \comp(U)=P1,\qquad P1=1Q,\qquad \comp(P)=Q.
\end{equation}
The first two tail stages are
\begin{equation}\label{eq:A-def}
 A=v_{N+1}=UUU11U,\qquad v_{N+2}=AAA11A.
\end{equation}

\begin{lemma}[complementary witness]\label{lem:witness}
The language of $X_N$ contains a word $W$ and its complement $\comp(W)$, both
of length
\[
 \ell_N=3L+1.
\]
\end{lemma}

\begin{proof}
At the outer spacer between the third and fourth copies of $A$ in
$v_{N+2}$, start at the second symbol of the final $U$ in the third $A$.
Because $U=0P$, the word of length $3L+1$ beginning there is
\begin{equation}\label{eq:W-def}
 W=P\,11\,U\,U.
\end{equation}
Indeed, the displayed pieces have total length
$(L-1)+2+L+L=3L+1$.

Now work inside the final copy of $A$ and start at its second $U$.  Reading
$3L+1$ symbols gives
\begin{equation}\label{eq:B-def}
 B=U\,U\,11\,Q,
\end{equation}
because only the length-$(L-1)$ prefix $Q$ of the final $U$ is used.  The
identities in \eqref{eq:PQ-identities} give
\begin{align*}
 B
 &= (Q0)(0P)11Q\\
 &=Q00(P1)(1Q)\\
 &=Q00(P1)(P1)\\
 &=\comp(P11UU)=\comp(W).
\end{align*}
Thus both $W$ and $\comp(W)$ occur in $v_{N+2}$ and hence in
$\Lang_{3L+1}(X_N)$.
\end{proof}

We next prove that this witness has the largest possible length.  A
\emph{defect} will mean the left endpoint of an equal adjacent pair.

\begin{lemma}[spacer and defect coding]\label{lem:defect-coding}
For $t\geq0$, put $T_t=v_{N+t}$.  There is a word
$S_t=a_1\cdots a_{4^t-1}\in\{0,1\}^{4^t-1}$ such that
\begin{equation}\label{eq:base-tiling}
 T_t=U(11)^{a_1}U(11)^{a_2}\cdots U(11)^{a_{4^t-1}}U.
\end{equation}
Here $(11)^0$ is empty.  Every finite boundary-type word in $V_N$ is a
factor of some $S_t$, and none contains $101$.

If $d_i$ is the defect at the boundary of type $a_i$, then its equal pair is
$00$ when $a_i=0$ and $11$ when $a_i=1$, and
\begin{equation}\label{eq:defect-distance}
 d_{i+1}-d_i=L+a_i+a_{i+1}.
\end{equation}
In particular consecutive defects have distance between $L$ and $L+2$.
\end{lemma}

\begin{proof}
The tail recursion gives $S_0=\eps$ and
\begin{equation}\label{eq:spacer-recursion}
 S_{t+1}=S_t0S_t0S_t1S_t,
 \qquad S_1=001.
\end{equation}
For $t\geq1$, induction shows that $S_t$ begins in $00$, ends in $01$, and
avoids $101$.  Indeed, a new junction has context either
$01\,0\,00=01000$ or $01\,1\,00=01100$, neither of which contains $101$.
Every finite boundary block of the nested limit occurs in some finite
$S_t$, so it has the same avoidance property.

The word $U$ is alternating, has odd length, and begins and ends in $0$.
Thus a type-$0$ boundary $UU$ contributes exactly the defect $00$, while a
type-$1$ boundary $U11U$ contributes exactly the middle defect $11$.
If $c_i$ is the start of the $U$ preceding boundary $a_i$, then
$d_i=c_i+L-1+a_i$ and $c_{i+1}=c_i+L+2a_i$.  Subtracting the corresponding
formulas for $d_i$ and $d_{i+1}$ proves \eqref{eq:defect-distance}.
\end{proof}

\begin{proposition}[exact complement separation]\label{prop:exact-separation}
The longest word $z$ for which both $z$ and $\comp(z)$ belong to the language
of $X_N$ has length $3L+1$.  Equivalently,
\begin{equation}\label{eq:exact-separation}
 s_{\Ftwo}(X_N)=3L+2=3\cdot2^{N+1}-1.
\end{equation}
\end{proposition}

\begin{proof}
The lower bound is \cref{lem:witness}.  For the upper bound, if
$z=z_0\cdots z_{q-1}$ is binary, set
\[
 E(z)=\{j\in\{0,\ldots,q-2\}:z_j=z_{j+1}\}.
\]
The set $E(z)$ and the first symbol determine $z$ recursively.  Hence
$E(z)$ alone determines exactly the pair $\{z,\comp(z)\}$.

Every language word can be placed with defects on both sides.  Namely, first
choose $t$ so that an occurrence is contained in $T_t$, then take the same
occurrence in the second copy of $T_t$ inside
$T_{t+1}=T_tT_tT_t11T_t$.  We use such interior occurrences for $z$ and
$\comp(z)$.  Their relative defect positions agree, and the boundary types
at corresponding visible defects are complementary.

Suppose first that at least three defects are visible.  For the first three,
write the two type triples as
$(a_i,a_{i+1},a_{i+2})$ and $(b_j,b_{j+1},b_{j+2})$.  Complementarity gives
$b_{j+r}=1-a_{i+r}$.  Equality of the two consecutive distances in
\eqref{eq:defect-distance} then gives
\[
 a_i+a_{i+1}=1,
 \qquad a_{i+1}+a_{i+2}=1.
\]
One triple is therefore $010$ and the other is $101$, contradicting
\cref{lem:defect-coding}.  Thus at most two defects are visible.

With no visible defect, the word lies between two consecutive defect starts,
so $q\leq L+2$.  With one visible defect, let $d_-<d<d_+$ be the preceding,
visible, and following starts.  The first coordinate of the word is at least
$d_-+1$, while its last coordinate is at most $d_+$.  Therefore
\begin{equation}\label{eq:one-defect-bound}
 q\leq d_+-d_-\leq2(L+2)\leq3L+1,
\end{equation}
where the last inequality uses $L\geq3$.

It remains to consider two visible defects.  Their equal spacing and
complementary types force the two type pairs to be $01$ and $10$; exchange
the occurrences so that the first has type $01$.  Align their first visible
defects at coordinate $0$.  Avoidance of $101$ forces the type preceding
$01$ to be $0$, so the preceding defect in the first occurrence starts at
$-L$.  The preceding defect in the $10$ occurrence starts at most at
$-(L+1)$.  Hence the common word starts no earlier than $-L+1$.

The second visible defect starts at $L+1$.  Avoidance of $101$ forces the
type following $10$ to be $0$, so the next defect in that occurrence starts
at $2L+1$; the next defect in the $01$ occurrence starts at least at
$2L+2$.  A defect whose left endpoint is the last coordinate of a word is
not internal, so the common word may end at $2L+1$ but no later.  Consequently
\[
 q\leq(2L+1)-(-L+1)+1=3L+1.
\]
This proves the upper bound, while \cref{lem:witness} attains it.  Restricting
that complementary pair also gives overlap at every shorter length, so the
first disjoint length is exactly $3L+2$.
\end{proof}

\begin{proposition}[exact inverse radius]\label{prop:radius-lower}
Let $R_N$ be the least $r\geq0$ for which the inverse of
$\Delta:X_N\to\Delta(X_N)$ has a local rule on $[-r,r]$.  Then
\[
 R_N=\frac{3L+1}{2}=3\cdot2^N-1.
\]
More generally, an inverse window $[-m,n]$ exists exactly when
$m+n\geq3L$.
\end{proposition}

\begin{proof}
By \cref{prop:exact-separation} and the exact window criterion
\eqref{eq:window-intro}, the condition is
$m+n+2\geq3L+2$, or $m+n\geq3L$.  The number
$L=2^{N+1}-1$ is odd.  Applying \eqref{eq:radius-intro},
\[
 R_N
 =\left\lceil\frac{3L}{2}\right\rceil
 =\frac{3L+1}{2}
 =\frac{3(2^{N+1}-1)+1}{2}
 =3\cdot2^N-1.
\]
\end{proof}

Combining
\cref{prop:minimality,prop:aperiodicity,prop:unique-ergodicity,cor:complement-disjoint,prop:radius-lower}
proves
\cref{thm:main-intro}.  The first few certified values are shown in
\cref{tab:values}; they are consequences of the formulas, not numerical
estimates.

\begin{table}[t]
\centering
\caption{Exact separation lengths and symmetric inverse radii.}
\label{tab:values}
\begin{tabular}{rrrr}
\toprule
$N$ & $|v_N|$ & $s_{\Ftwo}(X_N)$ & $R_N$\\
\midrule
1 & 3  & 11 & 5\\
2 & 7  & 23 & 11\\
3 & 15 & 47 & 23\\
4 & 31 & 95 & 47\\
5 & 63 & 191 & 95\\
\bottomrule
\end{tabular}
\end{table}

\begin{corollary}[no coarse uniform inverse bound]\label{cor:no-uniform}
There is no uniform bound on inverse symmetric radius depending only on
alphabet size $2$, cutting bound $4$, spacer bound $2$, and forward radius
$1$, even after restricting to minimal, aperiodic, uniquely ergodic rank-one
subshifts and to the fixed XOR derivative.
\end{corollary}

\begin{proof}
All four coarse parameters are independent of $N$, while
$R_N=3\cdot2^N-1\to\infty$.
\end{proof}
